\documentclass[../../main.tex]{subfiles}% 注意这里的文件路径不能用 ./main.tex ,否则用latexmk编译子文件会报错
\graphicspath{{\subfix{./image/}}{\subfix{../../image/}}} % 指定图片引用路径,后续可以直接使用图片文件名
% 如果图片存放在上级目录,则使用 ../../image/ 作为备用路径

% 例如:
% \begin{figure}[H]
% \centering
% \includegraphics[width=0.3\textwidth]{图.png}
% \caption{}
% \label{figure:图}
% \end{figure}
% 注意:上述\label{}一定要放在\caption{}之后,否则引用图片序号会只会显示??.

\begin{document}

\section{列紧性}

\begin{examples}
在度量空间$(C[0,1],\rho)$上，其中
\begin{align*}
\rho(x,y) \triangleq \max_{a \leqslant t \leqslant b} |x(t) - y(t)|,\quad \forall x,y\in C[0,1].
\end{align*}
考察点列
\begin{align*}
x_n(t)=
\begin{cases}
0, & t\geqslant\frac{1}{n}, \\
1-nt, & t<\frac{1}{n}
\end{cases}
\quad (n=1,2,\cdots).
\end{align*}
显然$\{x_n\}\subset B(\theta,1)$，其中$\theta$表示恒等于0的函数，但是$\{x_n\}$不含有$(C[0,1],\rho)$中的收敛子列。
\end{examples}
\begin{remark}
这个例子说明:在有穷维欧氏空间，有界无穷集必含有一个收敛子列，但这个性质不能推广到任意的度量空间。
\end{remark}

\begin{definition}
设$(\mathscr{X},\rho)$是一个度量空间，$A$为其一子集。称$A$是\textbf{列紧的}，如果$A$中的任意点列在$\mathscr{X}$中有一个收敛子列。若这个子列还收敛到$A$中的点，则称$A$是\textbf{自列紧的}。如果空间$\mathscr{X}$是列紧的，那么称$\mathscr{X}$为\textbf{列紧空间}。
\end{definition}

\begin{theorem}\label{theorem:列紧的基本性质}
\begin{enumerate}[(1)]
\item\label{theorem:列紧的基本性质-1} 在$\mathbb{R}^n$中任意有界集是列紧集，任意有界闭集是自列紧集。

\item\label{theorem:列紧的基本性质-2} 列紧空间内任意(闭)子集都是(自)列紧集。

\item\label{theorem:列紧的基本性质-3} 列紧空间必是完备空间。
\end{enumerate}
\end{theorem}
\begin{proof}
\begin{enumerate}[(1)]
\item 由\hyperref[Real Analysis-theorem:R^n中的Bolzano--Weierstrass定理]{Bolzano--Weierstrass定理}立得.

\item 定义自明.

\item 设$(\mathscr{X},\rho)$是一个列紧空间，$\{x_n\}$是其中的一串基本列，不妨设其是一无穷点列。由列紧性，存在$\{x_n\}$的子列$\{y_n\}$收敛到$x_0\in\mathscr{X}$。于是由\refpro{proposition:度量空间中基本列收敛当且仅当其子列收敛}可知$x_n\to x_0\ (n\to\infty)$.
\end{enumerate}
\end{proof}

\begin{definition}
设$M$是$(\mathscr{X},\rho)$度量空间中的一个子集，$\varepsilon>0$，$N\subseteq M$。如果对于$\forall x\in M,\exists y\in N$，使得$\rho(x,y)<\varepsilon$，那么称$N$是$M$的一个$\boldsymbol{\varepsilon }$\textbf{网}。如果$N$还是一个有穷集(个数依赖于$\varepsilon$)，那么称$N$为$M$的一个\textbf{有穷}$\boldsymbol{\varepsilon }$\textbf{网}。
\end{definition}

\begin{proposition}\label{proposition:泛函分析--网的等价定义}
设$M$是$(\mathscr{X},\rho)$度量空间中的一个子集，$\varepsilon>0$，$N\subseteq M$。则
\begin{enumerate}[(1)]
\item $N$是$M$的一个$\varepsilon$网当且仅当$M\subseteq\bigcup\limits_{y\in N}B(y,\varepsilon).$

\item $N$是$M$的一个有穷$\varepsilon$网当且仅当$M\subseteq\bigcup\limits_{y\in N}B(y,\varepsilon)$且$N$是一个有穷集(个数依赖于$\varepsilon$).
\end{enumerate}
\end{proposition}
\begin{proof}
定义自明.
\end{proof}

\begin{definition}
设$M$是$(\mathscr{X},\rho)$度量空间中的一个子集，$M$称为是\textbf{完全有界的}，如果$\forall\varepsilon>0$，都存在着$M$的一个有穷$\varepsilon$网。
\end{definition}

\begin{proposition}\label{proposition:完全有界必有界}
设$M$是$(\mathscr{X},\rho)$度量空间中的一个子集.
\begin{enumerate}[(1)]
\item\label{proposition:完全有界必有界-0} $M$是完全有界的当且仅当对$\forall \varepsilon >0$,存在有限个点$y_1,y_2,\cdots,y_n$,使得
\begin{align*}
M\subseteq \bigcup_{i=1}^n{B\left( y_i,\varepsilon \right)}.
\end{align*}

\item\label{proposition:完全有界必有界-3} $M$是完全有界的当且仅当$M$中任一序列都存在Cauchy子列.

\item\label{proposition:完全有界必有界-1} 若$M$是完全有界的,则$M$是有界集.

\item\label{proposition:完全有界必有界-2} 若$M$是有限集(有穷集),则$M$是完全有界的.
\end{enumerate}
\end{proposition}
\begin{proof}
\begin{enumerate}[(1)]
\item 由\refpro{proposition:泛函分析--网的等价定义}和完全有界的定义立得.

\item {\heiti 必要性:}假设 \(M\) 为完全有界集. 任取 \(M\) 中的一个序列 \(\{x_n\}_{n=1}^{\infty}\).
由完全有界的定义可知, 对任意 \(k \in \mathbb{N}\), 集合 \(M\) 能被有限个半径为 \(\frac{1}{k}\) 的开球覆盖.
因此, 存在一个半径为 \(1\) 的开球, 它包含 \(\{x_n\}\) 中的无穷多项. 记这些项的指标构成的集合为 \(S_1\).
同理, 存在一个半径为 \(\frac{1}{2}\) 的开球, 它包含 \(\{x_n\}_{n \in S_1}\) 中的无穷多项. 记这些项的指标构成的集合为 \(S_2\).
依此类推, 我们可以得到一列嵌套的指标集 \(S_1 \supseteq S_2 \supseteq \cdots \supseteq S_k \supseteq \cdots\), 且 \(S_k\) 中的元素对应的点都落在半径为 \(\frac{1}{k}\) 的同一个开球内.
选取严格递增的指标列 \(n_k \in S_k\), 构造 \(\{x_n\}\) 的一个子列 \(\{x_{n_k}\}\). 对任意 \(p,q \geqslant k\), 有 \(x_{n_p},x_{n_q} \in S_k\), 故它们位于同一个半径为 \(\frac{1}{k}\) 的开球中. 根据三角不等式, 有 \(\rho(x_{n_p},x_{n_q}) < \frac{2}{k}\). 由于 \(k\) 的任意性, 这表明 \(\{x_{n_k}\}\) 是 \(X\) 中的Cauchy列.

{\heiti 充分性:}假设 \(M\) 中任一序列都存在Cauchy子列. 用反证,若 \(M\) 不是完全有界的, 则存在 \(\varepsilon_0 > 0\), 使得 \(M\) 不能被有限个半径为 \(\varepsilon_0\) 的开球覆盖.
递归地构造 \(M\) 中的序列 \(\{x_n\}\).
任取 \(x_1 \in A\), 因 \(\{x_1\}\) 的 \(\varepsilon_0\) 邻域不能覆盖 \(M\), 必存在 \(x_2 \in A\) 使得 \(\rho(x_1,x_2) \geqslant \varepsilon_0\).
进一步地, 因 \(\{x_1,x_2\}\) 的 \(\varepsilon_0\) 邻域不能覆盖 \(M\), 必存在 \(x_3 \in A\) 使得 \(\rho(x_i,x_3) \geqslant \varepsilon_0\) (\(i = 1,2\)).
持续这个过程, 可构造出序列 \(\{x_n\}\) 使得对任意 \(m \neq n\), 均有 \(\rho(x_m,x_n) \geqslant \varepsilon_0\).
由于 \(\{x_n\}\) 的任意两项之间的距离都大于等于 \(\varepsilon_0\), 它的任意子列不可能成为Cauchy列, 这与题设矛盾!
因此 \(M\) 是完全有界集. 

\item 

\item 
\end{enumerate}
\end{proof}

\begin{proposition}[对角线法则]\label{proposition:对角线法则}
设 $\{a_{m,n}\}_{m,n=1}^{\infty}$ 是一族实数(或复数), 满足对每个固定的 $n$, 数列 $\{a_{m,n}\}_{m=1}^{\infty}$ 是有界的.
则存在严格递增的正整数序列 $\{m_k\}_{k=1}^{\infty}$, 使得对每个固定的 $n$, 子列 $\{a_{m_k,n}\}_{k=1}^{\infty}$ 都收敛.
\end{proposition}
\begin{proof}
对 $n=1$, 数列 $\{a_{m,1}\}_{m=1}^{\infty}$ 有界, 由 \hyperref[Real Analysis-theorem:R^n中的Bolzano--Weierstrass定理]{Bolzano--Weierstrass定理}, 存在收敛子列, 记其下标为
\[
m_1^{(1)} < m_2^{(1)} < m_3^{(1)} < \cdots,
\]
使得 $\{a_{m_k^{(1)},1}\}_{k=1}^{\infty}$ 收敛.

对 $n=2$, 考虑子列 $\{a_{m_k^{(1)},2}\}_{k=1}^{\infty}$, 它仍然有界, 因此又存在收敛子列, 记其下标为
\[
m_1^{(2)} < m_2^{(2)} < m_3^{(2)} < \cdots,
\]
且 $\{m_k^{(2)}\}_{k=1}^{\infty}$ 是 $\{m_k^{(1)}\}_{k=1}^{\infty}$ 的子列.
于是 $\{a_{m_k^{(2)},1}\}_{k=1}^{\infty}$ 和 $\{a_{m_k^{(2)},2}\}_{k=1}^{\infty}$ 都收敛.

如此继续, 对每个 $n$, 我们从上一行取出的子列中再取子列, 得到下标序列
\[
m_1^{(n)} < m_2^{(n)} < m_3^{(n)} < \cdots,
\]
满足 $\{m_k^{(n)}\}_{k=1}^{\infty}$ 是 $\{m_k^{(n-1)}\}_{k=1}^{\infty}$ 的子列, 并且对于所有 $j\leqslant n$, 数列 $\{a_{m_k^{(n)},j}\}_{k=1}^{\infty}$ 都收敛.最后得到一个无限的下标三角形矩阵
\begin{align*}
\begin{matrix}
m_{1}^{(1)}&		m_{2}^{(1)}&		m_{3}^{(1)}&		m_{4}^{(1)}&		...\\
m_{1}^{(2)}&		m_{2}^{(2)}&		m_{3}^{(2)}&		m_{4}^{(2)}&		...\\
m_{1}^{(3)}&		m_{2}^{(3)}&		m_{3}^{(3)}&		m_{4}^{(3)}&		...\\
\vdots&		\vdots&		\vdots&		\vdots&		\ddots\\
\end{matrix}
\end{align*}
现在定义最终的下标序列
\[
m_k \triangleq m_k^{(k)} \qquad (k=1,2,\cdots).
\]
则 $\{m_k\}_{k=1}^{\infty}$ 严格递增. 固定任意 $n\in \mathbb{N}$, 当 $k\geqslant n$ 时, 由于 $\{m_k^{(k)}\}_{k=n}^{\infty}$ 是 $\{m_k^{(n)}\}_{k=n}^{\infty}$ 的子列(因为每步取子列), 而 $\{a_{m_k^{(n)},n}\}_{k=1}^{\infty}$ 收敛, 所以其子列 $\{a_{m_k,n}\}_{k=n}^{\infty}$ 也收敛到同一极限. 因此 $\{a_{m_k,n}\}_{k=1}^{\infty}$ 收敛.
故对角线子列 $\{a_{m_k,n}\}_{k=1}^{\infty}$ 对每个固定的 $n$ 都收敛.
\end{proof}

\begin{theorem}[Hausdorff定理]\label{theorem:泛函分析--Hausdorff定理}
\begin{enumerate}[(1)]
\item 若度量空间$(\mathscr{X},\rho)$中的集合$M$是列紧的，则$M$是完全有界集。

\item 完备度量空间$(\mathscr{X},\rho)$中的集合$M$是列紧的当且仅当$M$是完全有界集。
\end{enumerate}
\end{theorem}
\begin{proof}
\begin{enumerate}[(1)]
\item 用反证法，若$\exists\varepsilon_0>0$，$M$中没有有穷的$\varepsilon_0$网。从而对$M$中任意有限个点$y_1,y_2,\cdots,y_n$,记$N=\{y_1,y_2,\cdots,y_n\}$,则有穷集$N$不是$M$的$\varepsilon_0$网,由\refpro{proposition:泛函分析--网的等价定义}知$M\nsubset \bigcup_{k=1}^n{B\left( y_k,\varepsilon _0 \right)}$,即$M\backslash \bigcup_{k=1}^n{B\left( y_k,\varepsilon _0 \right)}\neq \varnothing$.

任取$x_1\in M$，取$N=\{x_1\}$,则$\exists x_2\in M\setminus B(x_1,\varepsilon_0)$；

对$\{x_1,x_2\}\in M$，取$N=\{x_1,x_2\}$,则$\exists x_3\in M\setminus\left(B(x_1,\varepsilon_0)\cup B(x_2,\varepsilon_0)\right)$；

…………

对$\{x_1,x_2,\cdots,x_n\}\in M$，取$N=\{x_1,x_2,\cdots,x_n\}$,则$\exists x_{n+1}\in M\setminus\bigcup\limits_{k=1}^n B(x_k,\varepsilon_0)$；

…………

这样产生的点列$\{x_n\}\subset M$显然满足$\rho(x_n,x_m)\geqslant\varepsilon_0\ (n\neq m)$，故它的任意子列都不是基本列,因此由\refpro{proposition:收敛列必是基本列}知它没有收敛的子列。这与$M$的列紧性矛盾!

\item 必要性由(1)立得,下证充分性.对$M$中任意的无穷点列$\{z_n\}$，对$\forall \varepsilon >0$，都存在$M$的一个有穷$\varepsilon$网$N_{\varepsilon}$。
由命题知$M\subseteq \bigcup\limits_{y\in N_{\varepsilon}}B\left( y,\varepsilon \right)$。断言：存在$y_{\varepsilon}\in N_{\varepsilon}\subseteq M$，存在$\{z_n\}$的子列，使得$\{z_n^{(\varepsilon)}\}\subseteq B\left( y_{\varepsilon},\varepsilon \right)$，即$B\left( y_{\varepsilon},\varepsilon \right)$中包含$\{z_n\}$中无穷多项。

否则，对$\forall y\in N_{\varepsilon}$，$B\left( y,\varepsilon \right)$中只包含$\{z_n\}$的有限项。又$|N_{\varepsilon}|$有限，故$\bigcup\limits_{y\in N_{\varepsilon}}B\left( y,\varepsilon \right)$中至多包含$\{z_n\}$中有限项。因此$\{z_n\}\not\subseteq \bigcup\limits_{y\in N_{\varepsilon}}B\left( y,\varepsilon \right)$，这与$\{z_n\}\subseteq M\subseteq \bigcup\limits_{y\in N_{\varepsilon}}B\left( y,\varepsilon \right)$矛盾！

现在设$\{x_n\}$是$M$中的无穷点列.

取$\varepsilon =1$,$\{z_n\}=\{x_n\}$，则存在$y_1\in M$，存在$\{x_n\}$的子列$\{x_n^{(1)}\}$，使得$\{x_n^{(1)}\}\subseteq B\left( y_1,1 \right)$；

取$\varepsilon =\frac{1}{2}$,$\{z_n\}=\{x_n^{(1)}\}$，则存在$y_2\in M$，存在$\{x_n^{(1)}\}$的子列$\{x_n^{(2)}\}$，使得$\{x_n^{(2)}\}\subseteq B\left( y_2,\frac{1}{2} \right)$；

…………

取$\varepsilon =\frac{1}{k}$,$\{z_n\}=\{x_n^{(k-1)}\}$，则存在$y_k\in M$，存在$\{x_n^{(k-1)}\}$的子列$\{x_n^{(k)}\}$，使得$\{x_n^{(k)}\}\subseteq B\left( y_k,\frac{1}{k} \right)$；

…………

最后抽出对角线子列$\{x_k^{(k)}\}$。
事实上，$\forall\varepsilon>0$，当$n>\frac{2}{\varepsilon}$时，对$\forall p\in\mathbb{N}$有
\begin{align*}
\rho(x_{n+p}^{(n+p)},x_n^{(n)})\leqslant\rho(x_{n+p}^{(n+p)},y_n)+\rho(x_n^{(n)},y_n) 
\leqslant\frac{2}{n}<\varepsilon.
\end{align*}
故$\{x_k^{(k)}\}$是一个基本列。又因为$(\mathscr{X},\rho)$是完备的,所以$\{x_k^{(k)}\}$是$\{x_n\}$的一个收敛子列.因此$M$是列紧的.
\end{enumerate}
\end{proof}

\begin{definition}
一个度量空间若有可数的稠密子集，就称这个度量空间是\textbf{可分的}。
\end{definition}

\begin{theorem}
完全有界的度量空间是可分的。
\end{theorem}
\begin{proof}
设$(X,\rho)$是完全有界的度量空间,则对$\forall n\in \mathbb{N}$,取$N_n$为$X$的有穷的$\frac{1}{n}$网，则$\bigcup\limits_{n=1}^\infty N_n$是$X$的一个可数的稠密子集.事实上,显然$\bigcup\limits_{n=1}^\infty N_n$是可数的.对$\forall x\in X$,存在$y_k\in N_k(\forall k\in \mathbb{N})$,使得$\rho(x,y_k)<\frac{1}{k}$.进而$\lim_{k\to \infty}\rho(x,y_k)=0$.故$\bigcup\limits_{n=1}^\infty N_n$是$X$的一个稠密子集.
\end{proof}

\begin{definition}
在拓扑空间$\mathscr{X}$中，集合$M$称为是\textbf{紧的}，如果$\mathscr{X}$中每个覆盖$M$的开集族中有有穷个开集覆盖集合$M$.
\end{definition}

\begin{theorem}\label{theorem:紧集等价于自列紧}
设$(\mathscr{X},\rho)$是一个度量空间，为了$M\subseteq\mathscr{X}$是紧的当且仅当$M$是自列紧集.

进而由\hyperref[theorem:泛函分析--Hausdorff定理]{Hausdorff定理}知:若$M\subseteq\mathscr{X}$是紧的,则$M$是完全有界的.
\end{theorem}
\begin{remark}
若$(\mathscr{X},\rho)$是一个完备度量空间,$M$是完全有界的,则$M$不一定是紧的.此时由\hyperref[theorem:泛函分析--Hausdorff定理]{Hausdorff定理}只能得到$M$列紧,不一定自列紧!
\end{remark}
\begin{proof}

{\heiti 必要性:}设$M$是紧集。

(i)先证$M$是闭集，只要证$M$的余集是开集。$\forall x_0\in\mathscr{X}\setminus M$，因为
\begin{align*}
M\subseteq\bigcup\limits_{x\in M}B\left(x,\frac{1}{2}\rho(x,x_0)\right),
\end{align*}
利用$M$的紧性，$\exists x_k\in M\ (k=1,2,\cdots,n)$，使得
\begin{align*}
M\subseteq\bigcup\limits_{k=1}^n B\left(x_k,\frac{1}{2}\rho(x_k,x_0)\right).
\end{align*}
取$\delta=\min\limits_{1\leqslant k\leqslant n}\frac{1}{2}\rho(x_k,x_0)$，则显然有$\delta>0$，并且$\forall x\in B(x_0,\delta)$有
\begin{align*}
\rho (x,x_k)\geqslant \rho (x_k,x_0)-\rho (x_0,x)>\rho (x_k,x_0)-\delta \geqslant \frac{1}{2}\rho (x_k,x_0)\quad (k=1,2,\cdots ,n).
\end{align*}
因此$B(x_0,\delta )\cap \bigcup_{k=1}^n{B\left( x_k,\frac{1}{2}\rho (x_k,x_0) \right)}=\varnothing $，进而$B(x_0,\delta)\cap M=\varnothing$，即$B(x_0,\delta)\subseteq \mathscr{X}\setminus M$.从而$M$的余集是开集得证.

(ii)其次证$M$是列紧集。用反证法。假若有$M$中的点列$\{x_n\}$不含有收敛子列，不妨假定$x_n$是互异的。对每个$n\in\mathbb{N}$，做集合$S_n\triangleq\{x_1,x_2,\cdots,x_{n-1},x_{n+1},x_{n+2},\cdots\}$，显然每个$S_n$是闭集(因为不含收敛子列)，从而每个$\mathscr{X}\setminus S_n$是开集。但
\begin{align*}
\bigcup\limits_{n=1}^\infty(\mathscr{X}\setminus S_n)=\mathscr{X}\setminus\bigcap\limits_{n=1}^\infty S_n=\mathscr{X}\setminus\varnothing=\mathscr{X}\supseteq M,
\end{align*}
由$M$的紧性，$\exists N\in\mathbb{N}$，使得$\bigcup\limits_{n=1}^N(\mathscr{X}\setminus S_n)\supseteq M$，即得
\begin{align*}
\mathscr{X}\setminus\{x_n\}_{n=N+1}^\infty\supseteq M,
\end{align*}
但这是不可能的，因为$x_{N+1}$属于上式右端而不属于上式左端。此矛盾说明$M$是列紧的.又由(i)知$M$是闭集,所以其收敛子列都收敛到$M$中的点,故$M$是自列紧的.

{\heiti 充分性:}设$M$是自列紧的，要在$M$的任一开覆盖中取出有限覆盖。用反证法，如果某个开覆盖$\bigcup\limits_{\lambda\in\Lambda}G_\lambda\supseteq M$不能取出$M$的有限覆盖。由于$M$是自列紧的，$\forall n\in\mathbb{N}$，存在有穷的$\frac{1}{n}$网$N_n=\{x_1^{(n)},x_2^{(n)},\cdots,x_{k(n)}^{(n)}\}$，由\refpro{proposition:泛函分析--网的等价定义}知$\bigcup\limits_{y\in N_n}B(y,\frac{1}{n})\supseteq M$。因此，$\forall n\in\mathbb{N}$，$\exists y_n\in N_n$，使得$B(y_n,\frac{1}{n})$不能被有限个$G_\lambda$所覆盖。故这样取出来的点列$\{y_n\}$满足对$\forall n\in\mathbb{N}$,都有$B(y_n,\frac{1}{n})$不能被有限个$G_\lambda$所覆盖.
由假定$M$是自列紧集知，$\{y_n\}$必存在收敛子列$y_{n_k}$收敛到一点$y_0\in G_{\lambda_0}$。又$G_{\lambda_0}$是开集，所以$\exists\delta>0$，使得$B(y_0,\delta)\subseteq G_{\lambda_0}$。对此$\delta>0$，取$k$足够大，使得$n_k>\frac{2}{\delta}$，并且$\rho(y_{n_k},y_0)<\frac{\delta}{2}$，则$\forall x\in B(y_{n_k},\frac{1}{n_k})$，有
\begin{align*}
\rho(x,y_0)&\leqslant\rho(x,y_{n_k})+\rho(y_{n_k},y_0)\leqslant\frac{1}{n_k}+\frac{\delta}{2}<\delta,
\end{align*}
即$x\in B(y_0,\delta)$，从而$B(y_{n_k},\frac{1}{n_k})\subseteq B(y_0,\delta)\subseteq G_{\lambda_0}$。这与每个$B(y_n,\frac{1}{n})$不能被有限个$G_\lambda$所覆盖矛盾!
\end{proof}

\begin{theorem}\label{theorem:CMd是一个完备度量空间}
设$(M,\rho)$是一个紧的度量空间，用$C(M)$表示$M\to\mathbb{R}$的一切连续映射全体.定义
\begin{align*}
d(u,v)=\max\limits_{x\in M}|u(x)-v(x)|\quad (\forall u,v\in C(M)).
\end{align*}
则$(C(M),d)$是一个完备度量空间.
\end{theorem}
\begin{proof}
先证$(C(M),d)$是一个度量空间.
$d(u,v)$显然有非负性、唯一性、交换性和三角不等式.其实只有$d(u,v)$定义本身的合理性是要验证的。即对$\forall u\in C(M)$，存在着最大值$\max\limits_{x\in M}|u(x)|$。事实上，$u(M)$是紧集。因为对任意点列$y_n\in u(M)$，$\exists x_n\in M$，使得$u(x_n)=y_n$。由于$M$是紧的，从而有子列$x_{n_k}\to x_0,k\to\infty$，而$u$是连续的，便有$u(x_{n_k})\to u(x_0)\in u(M),k\to\infty$。令$y_0\triangleq u(x_0)$，即得$y_{n_k}=u(x_{n_k})\to y_0$，从而$u(M)$是紧集。由\hyperref[Real Analysis-theorem:Heine - Borel有限子覆盖定理]{Heine - Borel有限子覆盖定理}知$\mathbb{R}$中的紧集等价于有界闭集,故$u(M)\subseteq \mathbb{R}$是有界闭的数集。
由$\mathbb{R}$上的确界原理知，存在$\alpha=\inf u(M)$，$\beta=\sup u(M)$。再由$u(M)$是闭集知$\alpha,\beta\in u(M)$。故
\begin{align*}
\alpha=\min u(M),\quad \beta=\max u(M).
\end{align*}
这就证明了$\max\limits_{x\in M}|u(x)|$的存在性.

再证$(C(M),d)$是完备的.
设 $\{u_n\}_{n=1}^\infty$ 是 $C(M)$ 中的 Cauchy 列，即对任意 $\varepsilon>0$，存在 $N\in\mathbb{N}$ 使得当 $m,n\geqslant N$ 时，
\[
d(u_m,u_n)=\max_{x\in M}|u_m(x)-u_n(x)|<\varepsilon.
\]
于是对每个固定的 $x\in M$，当 $m,n\geqslant N$ 时，有
\begin{align}\label{eq::FFJIEJIOFJI234af}
|u_m(x)-u_n(x)|<\varepsilon.
\end{align}
即对每个固定的 $x\in M$,$\{u_n(x)\}$ 都是 $\mathbb{R}$ 中的 Cauchy 列，由 $\mathbb{R}$ 的完备性知它收敛，记
\[
u(x)=\lim_{n\to\infty}u_n(x),\quad \forall x\in M.
\]
这样就定义了函数 $u:M\to\mathbb{R}$。下面证明 $u\in C(M)$ 且 $d(u_n,u)\to0(n\to \infty)$。

首先证明$d(u_n,u)\to0(n\to \infty)$。对\eqref{eq::FFJIEJIOFJI234af}式令 $m\to\infty$，得
\begin{align}\label{eq::hiufh3FIOUEJO}
|u(x)-u_n(x)|\leqslant \varepsilon,\quad \forall x\in M,n\geqslant N.
\end{align}
因此对$\forall n\in \mathbb{N}$,有
\begin{align*}
d(u_n,u)=\max_{x\in M}|u(x)-u_n(x)|<\varepsilon.
\end{align*}
即$d(u_n,u)\to0(n\to \infty)$。

再证$u\in C(M)$。由\eqref{eq::hiufh3FIOUEJO}式知
\begin{align*}
|u_N(x)-u(x)|<\varepsilon,\quad \forall x\in M.
\end{align*}
而 $u_N$ 连续，故存在 $\delta>0$ 使当 $\rho(x,x_0)<\delta$ 时 
\begin{align*}
|u_N(x)-u_N(x_0)|<\varepsilon.
\end{align*}
于是
\[
|u(x)-u(x_0)|\leqslant |u(x)-u_N(x)|+|u_N(x)-u_N(x_0)|+|u_N(x_0)-u(x_0)|<3\varepsilon,
\]
所以 $u$ 在 $x_0$ 连续，从而 $u\in C(M)$。

因此每个 Cauchy 列都收敛于 $C(M)$ 中的元，即$(C(M),d)$ 完备.
\end{proof}

\begin{definition}
设$(M,\rho)$是一个紧的度量空间，用$C(M)$表示$M\to\mathbb{R}$的一切连续映射全体.
设$F$是$C(M)$的一个子集。称$F$是\textbf{一致有界的}，如果$\exists M_1>0$，使得$|\varphi(x)|\leqslant M_1(\forall x\in M,\forall\varphi\in F)$；

称$F$是\textbf{等度连续的}或\textbf{等度连续函数族}，如果$\forall\varepsilon>0$，总可以找到$\delta(\varepsilon)>0$，使得
\begin{align*}
|\varphi(x_1)-\varphi(x_2)|<\varepsilon\quad (\forall x_1,x_2\in M,\rho(x_1,x_2)<\delta,\forall\varphi\in F).
\end{align*}
\end{definition}

\begin{theorem}[Arzelà-Ascoli定理]\label{theorem:泛函分析--Arzelà-Ascoli定理}
设$(M,\rho)$是一个紧的度量空间，用$C(M)$表示$M\to\mathbb{R}$的一切连续映射全体.定义
\begin{align*}
d(u,v)=\max\limits_{x\in M}|u(x)-v(x)|\quad (\forall u,v\in C(M)).
\end{align*}
由\refthe{theorem:CMd是一个完备度量空间}知$(C(M),d)$是一个完备度量空间.
则$F\subseteq C(M)$是一个列紧集当且仅当$F$是一致有界且等度连续的函数族.
\end{theorem}
\begin{proof}
因为$C(M)$是完备的，所以由\hyperref[theorem:泛函分析--Hausdorff定理]{Hausdorff定理}知,$F$是列紧的当且仅当它是完全有界的.

{\heiti 必要性:}因为$F$是完全有界的，所以可设$N=\{\varphi_1,\cdots,\varphi_n\}$是$F$的有穷$1$网，则对$\forall \varphi \in F$，存在$\varphi_i\in N$，使得$d(\varphi,\varphi_i)<1$。于是
\begin{align*}
\max\limits_{x\in M}|\varphi(x)|\leqslant\max\limits_{x\in M}|\varphi(x)-\varphi_i(x)|+\max\limits_{x\in M}|\varphi_i(x)|=d(\varphi,\varphi_i)+\max\limits_{x\in M}|\varphi_i(x)|<1+\max\limits_{\substack{i=1,\cdots,n\\x\in M}} |\varphi_i(x)|.
\end{align*}
故对$\forall \varphi \in F,\forall x\in M$，有
\begin{align*}
|\varphi(x)|\leqslant 1+\max\limits_{\substack{i=1,\cdots,n\\x\in M}} |\varphi_i(x)|.
\end{align*}
因此$F$是一致有界函数族。

再证$F$是等度连续的.
对$\forall\varepsilon>0$，要证$\exists\delta=\delta(\varepsilon)$，使得$\forall\varphi\in F$有
\begin{align*}
|\varphi(x_1)-\varphi(x_2)|<\varepsilon\quad (\text{当}\ \rho(x_1,x_2)<\delta).
\end{align*}
因为$F$的$\frac{\varepsilon}{3}$网是一个有穷集$N(\frac{\varepsilon}{3})=\{\varphi_1,\varphi_2,\cdots,\varphi_n\}$，对这有穷个函数，由连续性，$\exists\delta=\delta(\frac{\varepsilon}{3})$，当$\rho(x_1,x_2)<\delta$有
\begin{align*}
|\varphi_i(x_1)-\varphi_i(x_2)|<\frac{\varepsilon}{3}\quad (i=1,2,\cdots,n).
\end{align*}
因为$\forall\varphi\in F,\exists\varphi_i\in N(\frac{\varepsilon}{3})$，使得$d(\varphi,\varphi_i)<\frac{\varepsilon}{3}$，所以
\begin{align*}
|\varphi(x)-\varphi(x')|&\leqslant|\varphi(x)-\varphi_i(x)|+|\varphi_i(x)-\varphi_i(x')|+|\varphi_i(x')-\varphi(x')| \\
&\leqslant 2d(\varphi,\varphi_i)+|\varphi_i(x)-\varphi_i(x')|<\varepsilon\quad (\text{当}\ \rho(x,x')<\delta).
\end{align*}

{\heiti 充分性:}设$F$一致有界且等度连续，对$\forall \varepsilon>0 $,我们要找有穷的$\varepsilon$网。由于$F$是等度连续的，$\exists\delta=\delta(\frac{\varepsilon}{3})>0$，使得当$\rho(x,x')<\delta$时，有
\begin{align}\label{eq::JIOFJ3423t58980}
|\varphi(x)-\varphi(x')|<\frac{\varepsilon}{3}\quad (\forall\varphi\in F).
\end{align}
由\refthe{theorem:紧集等价于自列紧}和\hyperref[theorem:泛函分析--Hausdorff定理]{Hausdorff定理}可知$M$是完全有界的,进而可取空间$M$上的有穷$\delta$网$N(\delta)=\{x_1,x_2,\cdots,x_n\}$。做映射
\begin{align*}
T:F\to\mathbb{R}^n,\quad T\varphi\triangleq(\varphi(x_1),\varphi(x_2),\cdots,\varphi(x_n))\quad (\forall\varphi\in F).
\end{align*}
记$\widetilde{F}=T(F)$，则$\widetilde{F}$是$\mathbb{R}^n$中的有界集。事实上，由$F$一致有界可设$|\varphi|\leqslant M_1\ (\forall\varphi\in F)$，则
\begin{align*}
\left( \sum_{i=1}^n{|\varphi (x_i)|^2} \right) ^{\frac{1}{2}}\leqslant \left[ \sum_{i=1}^n{\left( \max_{x\in M} |\varphi (x)| \right) ^2} \right] ^{\frac{1}{2}}\leqslant \sqrt{n}\max_{x\in M} |\varphi (x)|\leqslant \sqrt{n}M_1\quad (\forall \varphi \in F).
\end{align*}
从而$\widetilde{F}$是列紧集，利用\hyperref[theorem:泛函分析--Hausdorff定理]{Hausdorff定理}，$\widetilde{F}$有有穷的$\frac{\varepsilon}{3}$网$\widetilde{N}(\frac{\varepsilon}{3})=\{T\varphi_1,T\varphi_2,\cdots,T\varphi_m\}$。从而$\{\varphi_1,\varphi_2,\cdots,\varphi_m\}$是$F$的$\varepsilon$网，这是因为$\forall\varphi\in F,\exists\varphi_i$，使得
\begin{align}\label{eq::NFNAOIJOAA}
\rho _n(T\varphi ,T\varphi _i)=\left[ \left( \varphi \left( x_1 \right) -\varphi _i\left( x_1 \right) \right) ^2+\cdots +\left( \varphi \left( x_n \right) -\varphi _i\left( x_n \right) \right) ^2 \right] ^{\frac{1}{2}}=\left( \sum_{k=1}^n{\left[ \varphi \left( x_k \right) -\varphi _i\left( x_k \right) \right] ^2} \right) ^{\frac{1}{2}}<\frac{\varepsilon}{3}.
\end{align}
其中$\rho_n$表示$\mathbb{R}^n$上的欧氏距离.
于是对$\forall x\in M$,存在$x_r\in N(\delta)$，使得$\rho(x,x_r)<\delta$.再结合\eqref{eq::JIOFJ3423t58980}\eqref{eq::NFNAOIJOAA}式可得
\begin{align*}
|\varphi(x)-\varphi_i(x)|&\leqslant|\varphi(x)-\varphi(x_r)|+|\varphi(x_r)-\varphi_i(x_r)|+|\varphi_i(x_r)-\varphi_i(x)| \\
&<\frac{\varepsilon}{3}+\left( \sum_{k=1}^n{\left[ \varphi \left( x_k \right) -\varphi _r\left( x_k \right) \right] ^2} \right) ^{\frac{1}{2}}+\frac{\varepsilon}{3}\\
&<\frac{2}{3}\varepsilon+\rho_n(T\varphi,T\varphi_i)<\varepsilon,
\end{align*}
\end{proof}

\begin{example}
设$\Omega\subseteq\mathbb{R}^n$是有界开凸集。若$M_1,M_2$是两个给定的正数，则集合
\begin{align*}
F\triangleq\{\varphi\in C^{(1)}(\overline{\Omega})\mid |\varphi(x)|\leqslant M_1,|\mathrm{grad}\,\varphi(x)|\leqslant M_2(\forall x\in\Omega)\}
\end{align*}
是$C(\overline{\Omega})$上的一个列紧集，其中$C^{(1)}(\overline{\Omega})$表示$\overline{\Omega}$上的连续可微函数全体,即对$\forall \forall \varphi\in C^{(1)}(\overline{\Omega})$,有$\varphi$连续可微且$\varphi:\mathbb{R}^n\to \mathbb{R}$.
\end{example}
\begin{proof}
因为$\forall\varphi\in F,\forall x_1,x_2\in\overline{\Omega},\exists\theta\in(0,1)$，使得
\begin{align*}
\varphi(x_1)-\varphi(x_2)=\mathrm{grad}\,\varphi(\theta x_1+(1-\theta)x_2)\cdot(x_1-x_2),
\end{align*}
所以
\begin{align*}
|\varphi(x_1)-\varphi(x_2)|\leqslant M_2\rho_n(x_1,x_2)\quad (\forall\varphi\in F).
\end{align*}
其中$\rho_n$表示$\mathbb{R}^n$上的欧氏距离.这表明$F$是等度连续的。此外由$F$定义知$F$是一致有界的.再由\hyperref[theorem:泛函分析--Arzelà-Ascoli定理]{Arzelà-Ascoli定理}即得$F$是$C(\overline{\Omega})$上的一个列紧集.
\end{proof}

\begin{example}
在完备的度量空间$(X,\rho)$中求证: 子集 $A$ 列紧的充要条件是对 $\forall \varepsilon > 0$, 存在 $A$ 的列紧的 $\varepsilon$ 网.
\end{example}
\begin{proof}

{\heiti 必要性:} 由\hyperref[theorem:泛函分析--Hausdorff定理]{Hausdorff定理}知$A$列紧等价于$A$是完全有界的，故对$\forall \varepsilon >0$，存在$A$的有穷$\varepsilon$网$N_{\varepsilon}$。又由\hyperref[proposition:完全有界必有界-2]{有限集必是完全有界的}知$N_{\varepsilon}$是完全有界的，因此再由\hyperref[theorem:泛函分析--Hausdorff定理]{Hausdorff定理}知$N_{\varepsilon}$是$A$的列紧的$\varepsilon$网。

{\heiti 充分性:} 对$\forall \varepsilon >0$，设$N_{\varepsilon}$是$A$的列紧的$\frac{\varepsilon}{2}$网。由\hyperref[theorem:泛函分析--Hausdorff定理]{Hausdorff定理}知$N_{\varepsilon}$存在有穷的$\frac{\varepsilon}{2}$网$N_{\varepsilon}'$。
从而对$\forall a\in A$，存在$b\in N_{\varepsilon}$，使得$\rho(a,b)<\frac{\varepsilon}{2}$。进而存在$c\in N_{\varepsilon}'$，使得$\rho(b,c)<\frac{\varepsilon}{2}$。于是
\begin{align*}
\rho(a,c)\leqslant\rho(a,b)+\rho(b,c)<\varepsilon.
\end{align*}
故$N_{\varepsilon}'$也是$A$的有穷$\varepsilon$网，因此$A$是完全有界的。再由\hyperref[theorem:泛函分析--Hausdorff定理]{Hausdorff定理}知$A$是列紧的。
\end{proof}

\begin{proposition}\label{proposition:紧集上的连续函数必有最值}
设$(X,\rho)$是度量空间,则紧集$M\subseteq X$上的连续函数$f:M\to \mathbb{R}$必是有界的, 并且在$M$上达到它的上、下确界,即在$M$上存在最大、最小值.
\end{proposition}

\begin{example}
在度量空间中求证: 完全有界的集合是有界的, 并通过考虑 $l^2$ 的子集 $E = \{e_k\}_{k=1}^\infty$, 其中
\begin{align*}
e_k = \{\underbrace{0,0,\cdots,0}_{k-1},1,0,\cdots\},
\end{align*}
来说明一个集合可以是有界但不完全有界的.
\end{example}

\begin{example}
设 $(\mathscr{X},\rho)$ 是度量空间, $F_1,F_2$ 是它的两个紧子集, 求证: $\exists x_i \in F_i(i=1,2)$, 使得 $\rho(F_1,F_2) = \rho(x_1,x_2)$, 其中
\begin{align*}
\rho(F_1,F_2) \triangleq \inf\limits\{\rho(x,y) \mid x \in F_1, y \in F_2\}.
\end{align*}
\end{example}

\begin{example}
设 $M$ 是 $C[a,b]$ 中的有界集, 求证: 集合
\begin{align*}
\left\{F(x) = \int_a^x f(t)\mathrm{d}t \mid f \in M\right\}.
\end{align*}
是列紧集.
\end{example}

\begin{example}
设 $E = \{\sin nt\}_{n=1}^\infty$, 求证: $E$ 在 $C[0,\pi]$ 中不是列紧的.
\end{example}

\begin{example}
求证: \hyperref[proposition:空间S是完备度量空间]{$S$空间}的子集 $A$ 列紧的充要条件是: $\forall n \in \mathbb{N}, \exists C_n > 0$, 使得对 $\forall x = (\xi_1,\xi_2,\cdots,\xi_n,\cdots) \in A$, 有 $|\xi_n| \leqslant C_n\ (n=1,2,\cdots)$.
\end{example}

\begin{example}
设 $(\mathscr{X},\rho)$ 是度量空间, $M$ 是 $\mathscr{X}$ 中的列紧集, 映射 $f: \mathscr{X} \to M$ 满足
\begin{align*}
\rho(f(x_1),f(x_2)) < \rho(x_1,x_2) \quad (\forall x_1,x_2 \in \mathscr{X}, x_1 \neq x_2).
\end{align*}
求证: $f$ 在 $\mathscr{X}$ 中存在唯一的不动点.
\end{example}

\begin{example}
设 $(M,\rho)$ 是一个紧度量空间, 又 $E \subseteq C(M), E$ 中的函数一致有界并满足下列 Hölder 条件:
\begin{align*}
|x(t_1) - x(t_2)| \leqslant C\rho(t_1,t_2)^\alpha \quad (\forall x \in E, \forall t_1,t_2 \in M).
\end{align*}
其中 $0 < \alpha \leqslant 1, C > 0$. 求证: $E$ 在 $C(M)$ 中是列紧集.
\end{example}




\end{document}